\section{System Design}
\label{sec:design}

\system selects which Agents' histories deserve protection (O1),
which KV regions within those histories should be retained (O2), and
how much HBM replication may occupy as foreground demand changes (O3). 
Figure~\ref{fig:system-overview} shows our four cooperating components.
\emph{Tracker} estimates remaining workflow work and recovery urgency.
\emph{Router} chooses execution workers and replica destinations.
\emph{KV Master} selects Agent histories and KV regions, evaluates
candidate placements, and manages replica retention and reclamation.
\emph{Adapter} adjusts the HBM replication budget to foreground demand. Rank-local
Connectors execute transfers and report completion to KV Master.

\begin{figure}[t]
 \centering
 \includegraphics[width=\columnwidth]{figures/system-overview.pdf}
 \caption{\system: Four-component architecture. Each workflow node is one
 invocation. Letters identify Agents, and parentheses give invocation
 ordinals, B and C are each invoked twice.
 Each host contains one TP=2 DP worker. Each rank's HBM holds execution
 and replica regions. Rank-matched Connectors copy KV between hosts;
 each host's DRAM connects to its own Connectors. Blue paths carry
 KV; dashed paths carry control.}
 \label{fig:system-overview}
\end{figure}


\subsection{Tracker}
\label{sec:tracker}

% 中文：一个 Agent 通过多步完成其任务，每步对应一次模型调用。Tracker 按调用粒度
% 记录执行和历史使用关系，按 Agent 剩余工作的粒度估计其对工作流完成时间的影响。
An Agent completes its assigned task through a sequence of steps, each
corresponding to one model invocation. Tracker records execution and
history use at invocation granularity, and estimates remaining Agent
work to determine workflow urgency.

% 中文：本文考虑 Agent 任务 DAG 预先已知且固定的工作流。Agent 集合及依赖关系已知，
% 但每个 Agent 内部的调用次数和完成时间未知。
We target workflows with a fixed Agent-task DAG known before execution.
Agent identities and dependencies are fixed, while each Agent's
invocation count and completion time are unknown.

\subsubsection{Remaining-Time Prediction}
\label{sec:tracker-prediction}

% 中文：以 (w,a,r) 标识工作流 w 中 Agent a 的第 r 次调用，当依赖满足后，workflow runtime 将这次 invocation 提交给 serving 系统调度执行。
% 后续调用可以复用兼容前缀的已选择 KV；调用结束不等于历史失效，直到 Agent 结束。
Tracker identifies invocation $r$ of Agent $a$ in workflow $w$ as
$v=(w,a,r)$. Once its dependencies are satisfied, the workflow runtime
submits the invocation to the serving system for scheduling. Successive
invocations of the same Agent may reuse selected KV for a compatible
history prefix. Completing an invocation does not end the Agent's
history; Tracker tracks that history until the Agent finishes its task.

% 中文：每次调用提交时，Tracker 根据本次调用的信息预测 Agent 距离完成任务还需多久。
% 一次调用只预测一次，下一次调用到来时再更新，直到 Agent 结束。
At each invocation submission, Tracker predicts how long the Agent will
take to finish its task using the invocation's input and the Agent's
observed history. Each invocation triggers one prediction, which is
updated at each subsequent invocation until the Agent finishes its task.

% 中文：将预测建模为基于上下文的回归问题。文本编码器提取任务描述、本次新增输入
% 和近期已接受输出的语义表示，再与提交时已知的模型、角色、输入长度、已完成调用
% 次数、历史调用的长度和耗时，以及 Agent 已运行时间结合。
% 回归器直接预测 Agent 的剩余完成时间，供后续工作流 Slack 计算使用。
We formulate this prediction as context-conditioned regression. A text
encoder represents the Agent's task description, newly appended input,
and recent accepted outputs. A regressor combines this semantic
representation with features available at submission: the model and
Agent role, input length, completed-call count, previous calls' lengths
and timings, and elapsed Agent time. It directly predicts Agent $a$'s
remaining completion time, $\widehat R_{a,r}$, for workflow slack
computation.

% 中文：训练样本将每次调用提交时的特征与该次提交到 Agent 任务结束的实际时间配对。
Training examples pair these submission-time features with the observed
duration from submission of invocation $r$ to completion of Agent $a$'s
task.

% 中文：尚未启动的 Agent 使用同模型、同角色历史任务的时长先验 D。
% 首次调用提交时，以该先验初始化剩余时间估计 R，并记录提交时间 tau。
% 每次调用触发的预测返回后，以新预测替换 R，tau 始终取该预测对应的调用提交时间。
For an unstarted Agent, Tracker uses a duration prior $\bar D_a$ from
historical tasks with the same model and role. Its first invocation
submission initializes $\widehat R_a$ to this prior and $\tau_a$ to the
submission time. Each prediction $\widehat R_{a,r}$ replaces
$\widehat R_a$ when available, with $\tau_a$ set to invocation $r$'s
submission time.

\subsubsection{Slack Computation}
\label{sec:tracker-slack}

% 中文：注册工作流时，runtime 从固定编排图导出所有 Agent 任务的工作流 ID、Agent ID
% 和前驱 Agent ID 列表；边 a→b 表示 b 等待 a 完成任务后才能开始。
% runtime 报告任务完成事件，Tracker 移除已完成节点，得到剩余 DAG。
% 该 DAG 每个节点是一个尚未完成的 Agent 任务，其内部可以有多次 invocation。
At workflow registration, the runtime exports each Agent task's workflow
ID, Agent ID, and predecessor Agent IDs from the fixed orchestration graph.
An edge $a\to b$ means that $b$ waits for Agent $a$ to finish its task.
The runtime reports task completion, allowing Tracker to remove finished
tasks and maintain the residual DAG $G_t=(A_t,E_t)$. Each node in this
graph is one unfinished Agent task, which may contain multiple invocations.

% 中文：算法以剩余 Agent DAG、时长先验、已启动 Agent 的最新估计及提交时间、误差裕量为输入。
% 未启动 Agent 保留完整时长先验；已启动 Agent 的估计换算到当前时刻。
% 正向计算最早完成时间，反向计算最晚完成时间。
% 后向遍历以预测的工作流完成时间为终点；每个 Agent 的剩余工作只计算一次。
% pred 和 succ 分别表示剩余前驱与后继；g 是经校准的 Slack 高估误差裕量。
Algorithm~\ref{alg:slack-estimation} computes guarded Agent slack using
forward and backward critical-path passes~\cite{kelley1959cpm}.
Unstarted Agents retain their full duration priors, while estimates for
started Agents are adjusted to the current time $t$. The algorithm counts
each Agent's remaining work once and uses the predicted workflow finish
$T$ as the backward-pass endpoint. The two passes compute earliest and
latest finishes, $EF_a$ and $LF_a$. Here $\operatorname{pred}(a)$ and
$\operatorname{succ}(a)$ denote residual predecessors and successors,
and $g_a$ is a calibrated upper quantile of slack overestimation.

% 中文：所需时长估计缺失、过期或无效时，整个工作流暂不授予正 Slack。
% 正向遍历：EF = 本节点剩余时间 + 最晚完成的前驱；后向遍历：
% LF = 后继最晚开始时间的最小值，汇点从 T 开始回推。
% 最终 Slack 为 max(0, LF - EF - g)；两次遍历的复杂度均为 O(|A_t|+|E_t|)。
\par\medskip\noindent\begin{minipage}{\columnwidth}
\raggedright
\refstepcounter{designalgorithm}\label{alg:slack-estimation}
\textbf{Algorithm \thedesignalgorithm: Compute Agent slack}
\par\smallskip
\begin{tabularx}{\columnwidth}{@{}r@{\quad}X@{}}
\toprule
 & \textbf{Input:} $G_t=(A_t,E_t)$; Agent states; priors $\bar D_a$;
 estimates $(\widehat R_a,\tau_a)$; guards $g_a$; time $t$ \\
 & \textbf{Output:} guarded slack $\boldsymbol{s}=\{s_a:a\in A_t\}$ \\
\midrule
1 & $s_a\gets0,\quad \forall a\in A_t$ \\
2 & \textbf{for} $a\in A_t$ \textbf{do} \\
3 & \quad \textbf{if} $a$ is unstarted \textbf{then} $d_a\gets\bar D_a$ \\
4 & \quad \textbf{else} $d_a\gets\widehat R_a-(t-\tau_a)$ \\
5 & \textbf{if} any required estimate is missing or invalid, or
 any $d_a\le0$, \textbf{then return} $\boldsymbol{s}$ \\
6 & $\pi\gets\textsc{TopologicalOrder}(G_t)$ \\
7 & \textbf{for} $a$ in $\pi$ \textbf{do} \\
8 & \quad $EF_a\gets d_a+
 \max\bigl(\{0\}\cup\{EF_p:p\in\operatorname{pred}(a)\}\bigr)$ \\
9 & $T\gets\max\bigl(\{0\}\cup\{EF_a:a\in A_t\}\bigr)$ \\
10 & \textbf{for} $a$ in reverse $\pi$ \textbf{do} \\
11 & \quad $LF_a\gets
 \min\bigl(\{T\}\cup\{LF_b-d_b:b\in\operatorname{succ}(a)\}\bigr)$ \\
12 & \quad $s_a\gets\max\{0,LF_a-EF_a-g_a\}$ \\
13 & \textbf{return} $\boldsymbol{s}$ \\
\bottomrule
\end{tabularx}
\end{minipage}\par\medskip

% 中文：调用继承所属 Agent 的 Slack。
Invocation $v=(w,a,r)$ inherits its Agent's slack, $s_v=s_a$.

% 中文：预测和 Slack 更新在 Controller 中异步执行，调用调度和 token 生成无需等待。
Prediction and slack updates run asynchronously in the controller, so
invocation dispatch and token generation do not wait for them.

\subsection{Router and Connector}
\label{sec:router}

% 中文：Router 根据队列负载和前缀 KV 的位置调度调用，与此同时 Router 还会决定如果 Failure 发生时，该 Agent 调用 retry 到哪个DP Replica
% 这一选择也决定了如果要备份该调用对应的 KV，备份在哪里。
% retry 目标位于执行节点的受保护故障域之外，选择时考虑负载、可用内存和复制路径竞争。
Router schedules invocations based on queue load and prefix locality.
It also selects a retry DP replica for each Agent invocation in case of
failure. If the invocation's KV is selected for backup, it is placed at
this retry target. The target lies outside the execution worker's
protected failure domain and is chosen based on load, available memory,
and copy-path contention.

% 中文：已选 KV 存于 retry 目标的 GPU 内存（Peer HBM）或主机内存（Peer DRAM）。
% 配置的前 L 层中的已选区域按 KV Master 的优先级优先放入 HBM；溢出区域和后续层
% 的已选区域放入 DRAM。两层都是部分备份，每个已选区域只有一个稳定备份。
% 容量不足时，Router 将缺口返回 KV Master，重新选择保护区域。
At this target, selected KV resides in GPU memory (\emph{Peer HBM}) or
host memory (\emph{Peer DRAM}). Selected regions in the configured first
$L$ layers prefer Peer HBM according to KV Master's priorities; overflow
and selected regions in later layers use Peer DRAM. Both tiers hold
partial backups, with one stable backup per selected region.
If the selected set does not fit, Router returns the capacity shortfall
to KV Master for reselection.

% 中文：各 TP rank 的 Connector 执行 KV 复制并向 KV Master 报告完成状态。
% 故障后 Router 使旧执行失效，并将调用重定向到 retry 目标；Engine 直接复用目标
% HBM 中的 KV，Connector 将 DRAM 中的 KV 搬入 GPU。目标成为执行节点后，
% Router 为后续备份重新选择 Peer。
Rank-local Connectors execute KV copies and report completion to KV
Master. After a failure, Router fences the old execution and redirects
the invocation to the retry target. The engine reuses resident Peer HBM
KV, while Connectors move DRAM-resident KV into GPU memory.
If that target becomes the execution worker, Router selects a new peer
for subsequent backups.

% 中文：上述路由与传输机制处理已选状态；接下来由 KV Master 决定保护哪些
% Agent 历史，以及保留其中哪些 KV 区域。
These routing and transfer mechanisms operate on selected state.
KV Master determines which Agent histories and which KV regions to
protect, as described next.

\subsection{KV Master}
\label{sec:master}

% 中文：KV Master 是 \system 的策略核心，把 O1 和 O2 落实为两级选择：在 Agent 之间，
% 只把备份容量分配给恢复会推迟工作流的 Agent；在 Agent 内部，只备份缺失时会阻塞恢复的
% KV 区域。Tracker 提供 Slack，Router 提供 retry 目标及其 HBM、DRAM 容量。
KV Master is the policy core of \system. It realizes O1 and O2 as two
levels of selection. Across Agents, it spends backup capacity only where
recovery would delay the workflow. Within an Agent, it backs up only the
KV regions whose absence would stall recovery. Tracker supplies slack,
and Router supplies each history's retry target together with its HBM
and DRAM capacity.

% 中文：KV Master 不按 KV 大小、最近使用时间或重建代价排序，而是用“恢复延迟中暴露在
% 工作流上的部分”统一衡量所有决策。由此得到两个设计：一是 Slack 约束的保护，Slack 同时
% 决定是否保护一个 Agent 以及需要保护多少；二是感知暴露等待的区域选择，区域的价值取决于
% 它在逐层恢复中消除的等待，这也决定了它应放在 HBM 还是 DRAM。
Rather than ranking state by size, recency, or reconstruction cost, KV
Master measures every decision by the recovery delay it removes from the
workflow. Two design ideas follow. \emph{Slack-bounded protection} uses
an Agent's slack to decide both whether its history needs protection and
how much of it. \emph{Exposure-aware region selection} values a region by
the waiting it removes from layer-ordered recovery, which also
determines whether the region belongs in HBM or DRAM.

\subsubsection{Slack-Bounded Agent Selection}
\label{sec:agent-protection}

% 中文：设 S_a 为 Agent a 已选择备份的 KV 区域集合，T_a(S_a) 为其在 retry 目标上的
% 恢复延迟，即从中断的调用被重定向到恢复执行所需的时间；T_a(∅) 是没有备份时完整重建的
% 延迟。将 T_a(S_a) 和 Tracker 的 Slack 代入式 (1)，得到暴露在工作流上的延迟 E_a。
Let $S_a$ be the KV regions selected for Agent $a$, and let $T_a(S_a)$ be
its recovery delay on the retry target, from redirecting the interrupted
invocation to resuming its execution. $T_a(\emptyset)$ is the delay of
full reconstruction. Applying Equation~\ref{eq:isolated_penalty} with
$R_i=T_a(S_a)$ and Tracker's guarded slack $\sigma_i=s_a$ gives the
exposed delay
\begin{equation}
 E_a(S_a)=\bigl[T_a(S_a)-s_a\bigr]^+,
 \label{eq:exposed-delay}
\end{equation}
where $[x]^+=\max(0,x)$. Master selects state to reduce $E_a$ rather
than $T_a$.

% 中文：该度量有两个直接推论。第一，如果完整重建也能被 Slack 吸收，即 T_a(∅) ≤ s_a，
% 任何备份都无法缩短工作流，Master 不保护该 Agent。第二，对其余 Agent，收益在 T_a 降到
% s_a 时饱和，因此保护目标是未被吸收的延迟 τ_a，而不是完整检查点：Slack 充裕的 Agent
% 只需消除剩余的延迟，而关键路径上（s_a = 0）的 Agent 每缩短一秒恢复就缩短一秒工作流。
This metric has two consequences. First, if full reconstruction fits
within the slack, $T_a(\emptyset)\le s_a$, no backup can shorten the
workflow, and Master leaves the history unprotected. Second, for the
remaining Agents, the benefit saturates once $T_a(S_a)$ reaches $s_a$.
The protection target is therefore the unabsorbed delay
$\tau_a=T_a(\emptyset)-s_a$ rather than a complete checkpoint. An Agent
with ample slack needs only enough state to remove $\tau_a$, whereas
every second of recovery removed from a critical-path Agent, with
$s_a=0$, shortens the workflow by one second.

% 中文：图 2(b) 直观展示了这两点，图中的实心部分正是各 Agent 的保护目标。A1 的 33.2 s
% 重建完全落在 91.2 s 的 Slack 内，不分配备份容量；A2 没有 Slack，目标是完整的 32.6 s；
% A3 和 A4 只需分别消除 16.3 s 和 18.0 s。
Figure~\ref{fig:agent_value}(b) illustrates both consequences; its solid
segments are exactly the protection targets. A1's 33.2\,s reconstruction
fits within its 91.2\,s slack, so A1 receives no backup capacity. A2 has
no slack, so its target is the full 32.6\,s. A3 and A4 need to remove
only 16.3\,s and 18.0\,s, respectively.

% 中文：一次 worker 故障会中断其上的所有 Agent。由于 Tracker 在剩余 DAG 上计算 Slack，
% 同一汇合点上并列为瓶颈的前驱 Slack 都为 0，会同时获得高优先级；相对它们更早完成的
% 分支优先级较低。Tracker 没有有效估计时 Slack 为 0（算法 1），Master 保守地按恢复时间
% 的节省排序。计算 T_a(S_a) 需要知道恢复会用到哪些区域、缺少哪些区域会阻塞恢复，
% 这由下一节的区域选择回答。
A worker failure interrupts every Agent it hosts. Because Tracker
computes slack over the residual DAG, co-bottleneck predecessors of a
join all have zero slack and receive high priority together, whereas a
branch that finishes early relative to them is deprioritized. Without a
valid estimate, Algorithm~\ref{alg:slack-estimation} assigns zero slack,
and Master conservatively ranks the Agent by the recovery time it saves.
Computing $T_a(S_a)$ requires knowing which regions recovery will use
and which missing regions will stall it, which region selection
addresses next.

\subsubsection{Exposure-Aware Region Selection}
\label{sec:region-selection}

% 中文：KV 区域是某一层中的一个历史 token 块，包含其全部 TP 分片和所需的注意力元数据。
% 选择一个 Agent 并不会为其分配完整检查点，而是使其区域成为候选。
A KV region is one historical token block in one layer, including all its
TP shards and the attention metadata it requires. Selecting an Agent
makes its regions candidates rather than allocating a complete
checkpoint.

\paragraph{Estimating demand.}
% 中文：O2 表明，原生稀疏注意力只反复访问部分历史块。Master 利用引擎在正常执行中记录的
% 原生选择结果，估计 Agent a 第 ℓ 层区域 b 的访问因子。
O2 shows that native sparse attention repeatedly accesses a subset of
historical blocks. Master estimates this demand from the native
selections that the engine records during normal execution. For region
$b$ in layer $\ell$ of Agent $a$, the access factor is
\begin{equation}
 p_{a,\ell,b}=
 \frac{\text{selection units that selected }b}
      {\text{selection units for which }b\text{ was visible}}.
 \label{eq:demand}
\end{equation}
% 中文（续）：分母只统计因果上能看到该区域的选择单元，使较早和较新的区域可以比较。
% 统计遵循模型原生的选择粒度，跨 head 的重复选择只计一次。观测不足的区域使用校准的
% 逐层先验；稠密注意力中每个可见块的访问因子为 1。
Normalizing by visibility makes older and newer regions comparable.
Counts follow the model's native selection unit, counting duplicates
across heads once. Regions with few observations use a calibrated
per-layer prior, and every visible block under dense attention has
access factor one.

\paragraph{Modeling exposed waiting.}
% 中文：访问频率本身并不足够。接管后，被重定向的调用恢复执行，其前向计算逐层读取历史 KV，
% 因此 GPU 计算前面的层时，后面层的 KV 可以同时从 DRAM 传入。令 q_ℓ 为第 ℓ 层可以开始
% 消费历史的时间，α_ℓ 为其所需 KV 全部就绪的时间，c_ℓ 为其计算时间，其中包括按访问因子
% 加权的、对未备份区域的预期重建。
Access frequency alone is insufficient. After takeover, the redirected
invocation resumes with a forward pass that reads its history layer by
layer, so KV for later layers can be loaded from DRAM while the GPU
computes earlier layers. Let
$q_\ell$ be the time layer $\ell$ is ready to consume history,
$\alpha_\ell(S_a)$ the time its required KV becomes available, and
$c_\ell(S_a)$ its computation, including the expected reconstruction of
unbacked regions weighted by their access factors. Starting from
$q_1=0$,
\begin{equation}
 e_\ell(S_a)=\bigl[\alpha_\ell(S_a)-q_\ell\bigr]^+,\quad
 q_{\ell+1}=q_\ell+e_\ell(S_a)+c_\ell(S_a),
 \label{eq:overlap}
\end{equation}
% 中文（续）：e_ℓ 是计算无法掩盖的等待，T_a(S_a) 是最后一层完成的时间。Peer HBM 中的
% 区域在接管时即可使用；Peer DRAM 中的区域按层序经主机到设备链路传输，因此 α_ℓ 随前 ℓ
% 层所选 DRAM 字节数增长。重建计入未备份区域所依赖的低层状态，共享的重建工作只计一次。
where $e_\ell$ is the waiting that computation cannot hide, and
$T_a(S_a)$ is the finish time of the last layer. Regions in Peer HBM are
usable at takeover. Regions in Peer DRAM arrive in layer order over the
host-to-device link, so $\alpha_\ell$ grows with the DRAM bytes selected
for layers up to $\ell$. Reconstruction includes the lower-layer state
that unbacked regions depend on, counting shared work once.

\paragraph{Deriving tier placement.}
% 中文：式 (x) 区分了区域价值的两个来源。将区域备份到任一层，可避免按访问因子加权的重建；
% 放入 HBM 还能进一步消除其传输造成的暴露等待。暴露程度取决于某层之前有多少计算：第 1 层
% 之前没有计算，而更深层的传输可以与前面各层的计算重叠（如图 3(b) 中的 L2）。在各层字节
% 数相近时，暴露等待因此集中在前几层，这正是 Router 只允许前 L 层进入 Peer HBM 的原因。
% 若深层所选的 DRAM 字节超过前面计算所能掩盖的量，仍会产生等待（如图 3(b) 中的 L3）；
% 式 (x) 同样刻画这种情况，因此只有当避免的重建超过新增的等待时，Master 才备份这些区域。
% L 按模型和部署校准，取使后续层所选区域的传输能在这些层开始前完成的最小深度。
Equation~\ref{eq:overlap} separates two sources of region value. Backing
up a region in either tier avoids its demand-weighted reconstruction.
Placing it in HBM additionally removes any transfer waiting it would
expose. Exposure depends on how much computation precedes a layer.
Nothing precedes layer~1, whereas a deeper layer's transfer overlaps all
earlier layers, as for L2 in Figure~\ref{fig:kv_value}(b). With comparable
bytes per layer, exposed waiting therefore concentrates in early layers,
which is why Router admits only the first $L$ layers to Peer HBM. A deeper
layer can still stall when its DRAM transfer outlasts the preceding
computation, as for L3. Equation~\ref{eq:overlap} captures this case, so
Master backs up such regions only when the avoided reconstruction
outweighs the waiting they add. We calibrate $L$ per model and
deployment as the smallest depth for which profiled transfers of later
selected layers complete before those layers start.

% 中文：图 X 以 MiniCPM4-8B、32K token 历史为例量化这一效应，沿用第 2 节的就绪测量方式：
% retry 目标恢复历史 KV 并处理一个新增 token；KV 全部驻留时这一步需要 70 ms。四个面板
% 备份同一份部分状态：该 token 的原生选择实际读取的块，占 1 GiB 历史 KV 的 27%（272 MiB）；
% 未选的块在 HBM 和 DRAM 中都没有备份。(a) 单个恢复获得的有效 H2D 带宽为 6 GB/s 时，先从
% DRAM 传完备份再计算，暴露 47 ms 等待。(b) 按层序传输只剩第 1 层的 1.3 ms 等待，因为每层
% 备份平均 1.48 ms 即可传完，快于 2.18 ms 的计算。(c) 有效带宽降到 3 GB/s 时，每层传输需
% 2.97 ms，超过计算时间，等待累积到 27 ms，即图 3(b) 中 L3 的情形。(d) 将前 10 层的备份
% （86 MiB）放入 HBM 后，后续传输全部被计算掩盖，就绪时间回到 70 ms。作为对比，完整重建
% 需要 4.70 s 的 Prefill。把 272 MiB 备份全部放入 HBM 也不会更快，因此深度 L 以外的 HBM
% 不再缩短恢复。
Figure~\ref{fig:recovery-pipeline} quantifies this effect for MiniCPM4-8B
with a 32K-token history, following the readiness measurement of
Section~\ref{sec:motivation:cost}: the retry target restores the
historical KV and processes one additional token, which takes 70\,ms
with all KV resident. All panels back up the same partial state, namely
the blocks that this token's native selections read. These are 27\% of
blocks, or 272\,MiB of the 1\,GiB history, and unselected blocks have no
backup in either tier. When one recovery obtains an effective H2D
bandwidth of 6\,GB/s, loading the backup from DRAM before computation
exposes 47\,ms of waiting (a). Loading it in layer order hides all but
layer~1's 1.3\,ms (b), because a layer's backup loads in 1.48\,ms on
average, faster than its 2.18\,ms of computation. At 3\,GB/s, a layer's
backup takes 2.97\,ms to load, longer than its computation, and waiting
accumulates to 27\,ms (c), as for L3 in Figure~\ref{fig:kv_value}(b).
Serving the backup of the first ten layers, 86\,MiB, from HBM hides
every later transfer and restores readiness to 70\,ms (d). By
comparison, full reconstruction prefills the history in 4.70\,s. Placing
the entire 272\,MiB backup in HBM would be no faster, so HBM beyond
depth $L$ removes no further delay.

% 中文：因此 L 取决于单个恢复实际获得的有效 H2D 带宽。传输很少能打满链路，并发恢复还会
% 共享链路，因此这一带宽低于链路峰值。图 Y 对同一份备份扫描该带宽。高于 4.2 GB/s 时，传输
% 跟得上计算，只有第 1 层的传输会暴露，L = 1，只需 7.6 MiB HBM (a)；不使用 HBM 时的等待
% 也低于 2 ms (b)。低于这一点后，传输落后于计算，等待迅速增长，L 陡增：3 GB/s 时为 10，
% 2 GB/s 时为 17，1 GB/s 时为 25，此时不使用 HBM 的等待达到 217 ms。因为 L 只在这一点附近
% 及以下敏感，部署需要在预期的并发恢复下实测有效带宽来校准 L，而不是按链路标称峰值。
The depth therefore depends on the effective H2D bandwidth that one
recovery obtains. This bandwidth is below the link's peak, because
transfers rarely saturate the link and concurrent recoveries share it.
Figure~\ref{fig:depth-bandwidth} sweeps it for the same backup. Above
4.2\,GB/s, loads keep pace with computation, so only layer~1's transfer
is exposed and $L=1$, with 7.6\,MiB of HBM (a). Without HBM, the waiting
stays below 2\,ms (b). Below this point, loads fall behind computation,
waiting grows rapidly, and $L$ rises steeply: to 10 at 3\,GB/s, 17 at
2\,GB/s, and 25 at 1\,GB/s, where the waiting without HBM reaches
217\,ms. Because $L$ is sensitive only near and below this point, a
deployment calibrates it from the effective bandwidth measured under its
expected concurrent recoveries, not from the link's nominal peak.

% 中文（图示）：校准的解析模型，不是硬件实测：每层计算取 A100 上实测的单 token 前向时间，
% 有效带宽为假设值；备份是该 token 原生选择实际读取的块（事后统计），因此未备份的块在
% 就绪前不会被读取，重建不在时间线内。
\begin{figure}[t]
 \centering
 \includegraphics[width=\columnwidth]{figures/recovery-pipeline.pdf}
 \caption{Modeled readiness of one interrupted MiniCPM4-8B invocation
 with a 32K-token history: the retry target restores the historical KV
 and processes one additional token. All panels back up the same
 272\,MiB of the 1,024\,MiB history, namely the blocks that this token's
 native selections read, 27\% on average over three PG19 texts
 (retrospective). Unselected blocks have no backup. Each of the 32 layers
 computes in 2.18\,ms (measured on an A100) once its KV is available.
 The effective host-to-device bandwidth of the recovery is assumed to be
 6\,GB/s in (a, b) and 3\,GB/s in (c, d). (a) The backup loads from DRAM
 before computation. (b, c) It loads from DRAM in layer order. (d) The
 backup of layers 1--10 is in HBM, and the rest loads from DRAM. The
 model excludes failure detection, queueing, and interference.}
 \label{fig:recovery-pipeline}
 \Description{Four two-lane timelines, H2D above GPU, over time after
 takeover in milliseconds. Every panel backs up the same 272 of 1024 MiB
 of KV. (a) At 6 GB/s, the backup loads from DRAM for 47 ms while the GPU
 waits, then 32 layers compute; total 117 ms. (b) At 6 GB/s, the backup
 loads layer by layer and only layer 1 waits, for about 1 ms; total
 71 ms. (c) At 3 GB/s, loads take longer than computation, and waits
 accumulate between layers; total 97 ms with 27 ms of waiting. (d) At
 3 GB/s with layers 1 to 10 in HBM, later loads finish during computation
 and the GPU never waits; total 70 ms.}
\end{figure}

% 中文（图示）：同一解析模型与同一份备份，横轴为单个恢复获得的有效 H2D 带宽（对数刻度）。
% (a) 不产生等待所需的最小深度 L；(b) 备份全部在 DRAM、按层序传输时暴露的等待。虚线处
% （4.2 GB/s）以下传输落后于计算，L 超过 1；圆点为图 X 所用的两个带宽。
\begin{figure}[t]
 \centering
 \includegraphics[width=\columnwidth]{figures/depth-bandwidth.pdf}
 \caption{Required HBM depth versus the effective host-to-device
 bandwidth of one recovery, for the backup and readiness of
 Figure~\ref{fig:recovery-pipeline}. (a) Smallest depth $L$ for which no
 layer waits. (b) Exposed waiting when the whole backup is in DRAM and
 loads in layer order. Below the dashed line at 4.2\,GB/s, loads fall
 behind computation and $L$ exceeds one. Circles mark the two bandwidths
 of Figure~\ref{fig:recovery-pipeline}. Same analytical model.}
 \label{fig:depth-bandwidth}
 \Description{Two stacked panels sharing a logarithmic horizontal axis of
 effective H2D bandwidth from 0.25 to 32 GB/s. (a) The required depth L
 falls in steps from 31 at 0.25 GB/s to 10 at 3 GB/s and reaches 1 at
 4.2 GB/s, where a dashed line marks where loads keep pace with
 computation; it stays 1 above. (b) On a logarithmic scale, the waiting
 without HBM falls from about 1000 ms at 0.25 GB/s to 27 ms at 3 GB/s,
 drops sharply near 4.2 GB/s, and declines from about 2 ms to 0.25 ms at
 32 GB/s. Circles mark 3 and 6 GB/s.}
\end{figure}

\paragraph{Selecting regions.}
% 中文：对一个受保护的 Agent，Master 按访问因子从高到低加入区域，每加入一个就用式 (x)
% 重新计算 E_a，直到恢复延迟落入 Slack，或继续加入不再降低 E_a。对前 L 层中已选的区域，
% Master 按“放入 HBM 比放入 DRAM 每字节多消除的暴露延迟”向 Router 提供优先级；HBM 放不下
% 的区域退回 DRAM，Master 将这些延后的 HBM 候选报告给 Adapter，作为增大配额的信号。
For a protected Agent, Master adds regions in decreasing access factor
and re-evaluates $E_a$ with Equation~\ref{eq:overlap} after each
addition. It stops when the recovery delay fits within the slack or
further regions no longer reduce $E_a$. For selected regions in the
first $L$ layers, Master gives Router priorities based on the exposed
delay that HBM avoids over DRAM per byte. Regions that do not fit in HBM
fall back to DRAM, and Master reports these deferred HBM candidates to
Adapter as its signal for growing the quota (Section~\ref{sec:capacity}).

% 中文：图 X 以一个受保护的 Agent 为例说明这一过程。(a) 假设访问因子最高的 5 个区域
% 足以达到保护目标。(b) L = 2、HBM 配额为 2 个区域时，Router 将 L1/B0 和 L1/B2 放入
% HBM，因为第 1 层之前没有计算可以掩盖其传输；L2/B1 因配额已满退回 DRAM，其传输可与第 1
% 层的计算重叠；第 3 层的已选区域超出深度 L，放入 DRAM。
Figure~\ref{fig:selective-recovery} walks through one protected Agent.
Suppose its five hottest regions meet the protection target (a). With
$L=2$ and an HBM quota of two regions, Router places L1/B0 and L1/B2 in
HBM, because no computation precedes layer~1 to hide their transfer.
L2/B1 overflows to DRAM, where its load overlaps layer-1 computation,
and the selected layer-3 regions lie below depth $L$ and go to DRAM (b).

% 中文（图示）：构造示例，访问因子均为示意，不是测量结果。
\begin{figure}[t]
 \centering
 \includegraphics[width=\columnwidth]{figures/selective-recovery.pdf}
 \caption{Region selection and placement for one protected Agent
 (constructed example). (a) Access factors per layer and block. Master
 adds regions in decreasing $p$; the five hottest (solid) are assumed to
 meet the protection target. (b) With $L=2$ (dashed line) and an HBM
 quota of two regions, L1/B0 and L1/B2 occupy the retry target's HBM,
 since no computation precedes layer~1 to hide their transfer. L2/B1
 overflows to its DRAM, and selected layer-3 regions use DRAM.
 Unselected regions (--) have no backup.}
 \label{fig:selective-recovery}
 \Description{Two three-by-four grids of layers L1 to L3 and blocks B0 to
 B3. The first grid lists access factors; the five regions with the
 highest values, L1/B0, L1/B2, L2/B1, L3/B0 and L3/B2, are selected. The
 second grid shows L1/B0 and L1/B2 in HBM and the other three selected
 regions in DRAM, with a dashed line below layer 2 marking depth L.}
\end{figure}

\paragraph{Correctness.}
% 中文：副本只影响恢复时间，不影响正确性：任何未备份的区域都可以由持久化的 token 历史
% 重建。只有当引擎能够将保留的区域与缺失区域的重建组合使用时，Master 才计入部分备份的
% 收益；否则 T_a 按完整前缀 Prefill 计算，部分副本不获得避免计算的收益。
Replicas affect recovery time but never correctness, because any unbacked
region can be reconstructed from the durable token history. Master
credits partial restoration only when the engine can combine retained
regions with reconstruction of the missing ones. Otherwise, $T_a$ models
full-prefix prefill, and partial replicas receive no credit for avoided
computation.

\subsubsection{Replica Management}
\label{sec:master-online}
\label{sec:execution}

% 中文：选择随 Agent 执行持续更新。调用提交与完成、Tracker 的 Slack 更新以及 Adapter 的
% 配额变化都会触发重新规划。历史以追加方式增长，Master 在新的 KV 块写满时将其加入候选，
% 只复制新选择或位置改变的区域。后续调用扩展同一前缀，因此已有副本跨调用保持有效。
Selections evolve as Agents execute. Invocation submissions and
completions, slack updates, and quota changes trigger replanning.
Because histories grow by appending, Master adds each KV block to the
candidates once it fills and copies only regions that are newly selected
or moved. Existing replicas
remain valid across invocations, since later invocations of the same
Agent extend the same prefix.

% 中文：在 Agent 之间，Master 以“每字节消除的暴露延迟”作为统一尺度分配容量。算法 2 把
% 所有受保护 Agent 的候选区域放入同一个优先队列并贪心接纳。关键路径上的 Agent 每消除一秒
% 恢复就缩短一秒工作流，因此自然优先；已达到保护目标的 Agent，其后续区域收益为 0，不再
% 占用容量。由于区域收益会随同一 Agent 其他区域的接纳而变化，Master 会重新评估弹出的
% 候选，若其收益已低于队列中下一个估计值，则重新入队。
Across Agents, Master allocates capacity by exposed delay removed per
byte, a common scale for Agents with different slack.
Algorithm~\ref{alg:protection} places the candidate regions of all
protected Agents in one priority queue $Q$ and admits them greedily.
Critical-path Agents shorten the workflow by every second of recovery
removed and therefore come first, whereas an Agent that has met its
target gains nothing from further regions and stops consuming capacity.
Because a region's gain changes as other regions of its Agent are
admitted, Master re-evaluates each popped candidate under Router's
placement and reinserts it if its gain has fallen below the best
remaining estimate in $Q$.

% 中文（算法）：先根据完整重建延迟与 Slack 判断哪些 Agent 需要保护，只把它们的候选区域
% 放入队列；然后按单位字节收益贪心接纳，已达目标的 Agent 跳过，收益下降的候选重新入队；
% 容量不足时，只有收益高出 ε 才替换已有副本；最后复制并发布新区域，在读取结束后释放被
% 移除的区域。
\par\medskip\noindent\begin{minipage}{\columnwidth}
\raggedright
\refstepcounter{designalgorithm}\label{alg:protection}
\textbf{Algorithm \thedesignalgorithm: Slack-bounded replica selection}
\par\smallskip
\begin{tabularx}{\columnwidth}{@{}r@{\quad}X@{}}
\toprule
 & \textbf{Input:} slack $s_a$; access factors $p$; selections
 $\{S_a\}$; Router targets and capacities \\
 & \textbf{Output:} updated selections and placements \\
\midrule
1 & \textbf{for} each Agent $a$ with a materialized history \textbf{do} \\
2 & \quad \textbf{if} $T_a(\emptyset)>s_a$ \textbf{then} add $a$'s unselected
 regions to $Q$ \\
3 & \textbf{while} $Q\neq\emptyset$ and planning budget remains \textbf{do} \\
4 & \quad pop $x$ of Agent $a$ with the largest estimated gain \\
5 & \quad \textbf{if} $E_a(S_a)=0$ \textbf{then continue}
 \hfill$\triangleright$ target met \\
6 & \quad $g\gets\bigl[E_a(S_a)-E_a(S_a\cup\{x\})\bigr]/\mathrm{bytes}(x)$ \\
7 & \quad \textbf{if} $g<\textsc{TopKey}(Q)$ \textbf{then} reinsert $x$ with key $g$ \\
8 & \quad \textbf{else if} $x$ fits \textbf{then} $S_a\gets S_a\cup\{x\}$ \\
9 & \quad \textbf{else} evict victims with gain $<g-\epsilon$ to fit $x$ \\
10 & Copy added regions; publish them when complete \\
11 & Release removed regions after active reads finish \\
\bottomrule
\end{tabularx}
\end{minipage}\par\medskip

% 中文：容量不足时，新候选可以替换单位字节收益更低的已有副本，但只有收益高出至少 ε 时
% 才替换，以免副本在收益相近的候选之间反复迁移。Slack 增大到足以吸收恢复的 Agent，其
% 副本收益为 0，最先被替换。
When capacity is short, a candidate may replace replicas with lower gain
per byte, but only if its gain exceeds theirs by a margin $\epsilon$,
which prevents replicas from oscillating between similar candidates.
Replicas of Agents whose slack has grown to absorb their recovery have
zero gain and are replaced first.

% 中文：副本的生命周期跨越单次调用：一次调用结束不会回收其 Agent 的副本，因为下一次调用
% 依赖同一历史。Agent 完成任务后，Master 释放其引用，副本在正在进行的读取和传输结束后回收。
Replicas outlive individual invocations. Completing an invocation does
not retire its Agent's replicas, because the next invocation depends on
the same history. When the Agent finishes its task, Master releases its
references, and the replicas are reclaimed once active reads and
transfers finish.

% 中文：复制是异步的，因此 Master 用副本目录保证恢复只使用完整的数据。目录记录每个区域的
% 层、token 区间、TP 分片、兼容性元数据和位置；Connector 报告复制完成，区域只有在所有分片
% 和元数据都到达后才发布。迁移期间新旧两份副本都计入容量，直到目的副本发布且旧副本上的
% 读取结束。
Because copies are asynchronous, a replica manifest ensures that
recovery uses only complete data. The manifest records each region's
layer, token range, TP shards, compatibility metadata, and location.
Connectors report copy completion, and a region is published only after
all its shards and metadata arrive. During a move, both copies count
toward capacity until the destination is published and reads of the old
copy finish.

% 中文：故障后，Router 将中断的调用重定向到 retry 目标。引擎直接使用 HBM 中的区域，
% Connector 按层序把 DRAM 中的区域传入 GPU，引擎按需从持久化的输入重建缺失的区域，
% 包括尚未完成复制的最新 token 的 KV。目录显式记录未覆盖的区域，因此部分副本不会被当作
% 完整前缀使用。
After a failure, Router redirects the interrupted invocation to its
retry target. The engine uses HBM regions in place, Connectors load DRAM
regions in layer order, and the engine reconstructs missing regions on
demand from the durable inputs. As in checkpoint-and-replay
recovery~\cite{strati2024dejavu,cao2026lumen}, this includes KV for
accepted tokens whose copies had not completed.
Because the manifest records coverage gaps explicitly, partial replicas
are never treated as a complete prefix.

\subsection{Adapter}
\label{sec:capacity}

Adapter realizes O3 by determining the HBM replication capacity
foreground execution can support. For takeover rank $k$,
$B^{\rm pool}_k$ excludes weights and fixed allocations,
$B^{\rm serve}_k$ covers the serving high watermark and execution
reservations, and $B^{\rm reserve}_k$ reserves growth and workspace.
The ceiling is
\begin{equation}
 C_k=\max\{0,\min(B^{\rm cap}_k,
 B^{\rm pool}_k-B^{\rm serve}_k-B^{\rm reserve}_k)\}.
 \label{eq:ceiling}
\end{equation}
HBM replica use is not part of $B^{\rm serve}_k$ and is not subtracted twice.
Committed, reserved, copying, and retiring allocations all count toward
physical occupancy. Distinct TP ranks have distinct ceilings.

Adapter uses additive increase and multiplicative
decrease~\cite{chiu1989aimd}. Allocation stalls, reserve violations, or
memory-related admission waits multiply quota $B_k$ by $0<\gamma<1$.
After $m$ stable windows below a lower wait threshold, Adapter adds
$\Delta$ bytes if Master reports beneficial early-layer candidates
currently in Peer DRAM or deferred for HBM capacity. The signal is computed
before quota filtering, allowing growth from zero. Updates are
page-aligned and clipped to $C_k$; missing telemetry disables growth.
Adapter changes total quota, while the configured depth $L$ remains
unchanged.

When quota shrinks, Master stops new HBM admissions and asks Router
to move the least valuable HBM placements to the same peer's DRAM if
their regions remain selected. The plan reserves that host's DRAM
space; Connectors copy and validate the destination. Master then
redirects the manifest and releases the HBM allocation after readers
drain. If that peer's DRAM cannot accommodate the migration, Master
reselects the finite backup set, explicitly removing regions and their
predicted protection benefit as needed.
Their token histories still permit reconstruction. Mandatory pressure
relief may increase exposed delay, but cannot bypass active physical leases.
Foreground admission uses actual occupancy; a reduced quota alone
does not free memory. Pending copies may be cancelled, and allocations
wait when immediate physical space is unavailable.

The four components specify a design. Integrated measurements are
still needed to establish recovery gains, cross-invocation reuse, and
normal-service overhead.
